\origpage{45}{N. Robertson, P.D. Seymour / Journal of Combinatorial Theory, Series B 89 (2003) 43--76}
\noindent graph. It can be shown that there is no class of graphs of bounded genus such that all these graphs are tree-structures over it, and that shows the necessity for our third ingredient. Let $\Sigma$ be a surface, let $N \ge 0$ be an integer, and let $L$ be a graph such that no graph obtained from $L$ by deleting $\le N$ vertices can be drawn in $\Sigma$. Then there will be no $L$ minor in any graph which can be constructed by adding $\le N$ vertices (joined arbitrarily) to a graph drawn in $\Sigma$. This then yields a third type of structure which we should anticipate. If a subgraph $G'$ of $G$ can be obtained from $G$ by deleting $\le n$ vertices of $G$, let us call $G$ a \textit{($\le n$)-vertex extension} of $G'$.

We conjectured for some time that these three were sufficient. But that was false.

\textit{Ingredient} 4: \textit{Vortices.} Take a graph drawn in the plane, and let the vertices on the infinite region be $v_1, v_2, \dots, v_n$ in order. Add new edges joining $v_1$ to $v_3, v_2$ to $v_4, v_3$ to $v_5$ and so on, and let the resulting graph be $G$. Then it can be shown that such a graph has no large clique minor; indeed, Seese and Wessel [11] found the exact bound, that $G$ can have a $K_7$, but cannot have a $K_8$, minor. Yet such graphs $G$ cannot be constructed from our first three ingredients, and so a fourth ingredient is needed, so-called ``vortices''. Roughly, a vortex is a graph with some of its vertices arranged in a circular order, so that however these special vertices are partitioned into two intervals, there are only a bounded number of disjoint paths from one interval to the other. For instance, the graph $G$ above consists of a planar graph together with a vortex (formed by the new edges) inserted into the infinite region.

Our main result implies that these four ingredients suffice. Let us turn to the statement of the theorem. We need first to define what we mean by an ``$r$-ring with perimeter $t_1, \dots, t_n$''. Roughly speaking, it consists a graph $G$ and a sequence $t_1, \dots, t_n$ of distinct vertices of $G$, such that $G$ can be constructed as follows. Let $G_1, \dots, G_n$ be mutually disjoint graphs, each with $\le r$ vertices, and for $1 \le i \le n$ let $t_i \in V(G_i)$. Let us choose some $i$ with $1 \le i \le n$, choose $u \in V(G_i)$ and $v \in V(G_{i+1})$ and identify $u$ with $v$, provided that $t_1, \dots, t_n$ remain all distinct; and repeat this process as often as we wish.

More precisely, we say $G$ is an \textit{$r$-ring with perimeter $t_1, \dots, t_n$} if $t_1, \dots, t_n \in V(G)$ are distinct and there is a sequence $X_1, \dots, X_n$ of subsets of $V(G)$, such that

\hspace*{12pt}$\bullet$ $X_1 \cup \dots \cup X_n = V(G)$, and every edge of $G$ has both ends in some $X_i$,

\hspace*{12pt}$\bullet$ $t_i \in X_{\delm{1}\fix{i}}$ for $1 \le i \le n$,\fixnote{fix:r-ring}{As printed, all of $t_1, \dots, t_n$ lie in $X_1$, which has at most $r$ vertices, so an $r$-ring could have at most $r$ perimeter vertices. The sets $X_1, \dots, X_n$ are a path-decomposition laid along the perimeter, one for each $t_i$.}

\hspace*{12pt}$\bullet$ $X_i \cap X_k \subseteq X_j$ for $1 \le i \le j \le k \le n$,

\hspace*{12pt}$\bullet$ $|X_i| \le r$ for $1 \le i \le n$.

It is easy to check that this is equivalent with the previous definition.

Now let $G_0$ be a graph drawn in a surface $\Sigma$, and let $\Delta_1, \dots, \Delta_d \subseteq \Sigma$ be pairwise disjoint closed discs, each meeting the drawing only in vertices of $G_0$, and each containing no vertices of $G_0$ in its interior. For $1 \le i \le d$ let the vertices of $G_0$ in $\text{bd}(\Delta_i)$ be $t_1, \dots, t_n$ say, in order, and choose an $r$-ring $G_i$ with perimeter $t_1, \dots, t_n$, meeting $G_0$ just in $t_1, \dots, t_n$ and disjoint from every other $G_j$; and let $G$ be the union of $G_0, G_1, \dots, G_d$. Such a graph $G$ (and any graph isomorphic to it) is called an \textit{outgrowth by $d$ $r$-rings of a graph in $\Sigma$}. Now we can state our theorem.
